\section{Primene u telekomunikacijama}
\label{sec:telekom}
% TODO: Razraditi poglavlje tek nakon provere izvora; ne unositi neproverene činjenice.
Analiza primena pratiće tri nivoa: AES kao primitivu, način rada i protokol
koji određuje format poruke i upravljanje ključevima. Početne tačke su
specifikacije IPsec ESP sa GCM i istorijska specifikacija TLS 1.3
iz 2018.~\cite{rfc4106,rfc8446}; podrška algoritmu ne dokazuje obim njegove
stvarne upotrebe.

\subsection{Mrežni slojevi i zaštita podataka u tranzitu}
% TODO: Razdvojiti pristupnu, transportnu i aplikativnu zaštitu, granice poverenja i upravljanje ključevima.

\subsection{VPN i IPsec}
% TODO: Obraditi ESP, AES-GCM i istorijske režime; odvojiti IKE razmenu ključeva od šifrovanja saobraćaja, uz tačne RFC odeljke.

\subsection{TLS i transportni protokoli}
% TODO: Mapirati konkretne verzije i cipher suites na primitive i AEAD; proveriti naslednike, errata i važenje RFC 8446 na datum preseka.

\subsection{Wi-Fi, WPA2 i WPA3}
% TODO: Proveriti IEEE 802.11 i Wi-Fi Alliance dokumente za CCMP/GCMP i profile; razlikovati autentifikaciju SAE od šifrovanja saobraćaja. Ne tvrditi da svaki profil koristi isti mod ili ključ.

\subsection{Mobilni sistemi}
% TODO: Istražiti tačne 3GPP specifikacije i release za AES-zasnovane EEA2/NEA2 i EIA2/NIA2 gde je relevantno; odvojiti enkripciju od integriteta i od drugih algoritama.

\subsection{Infrastruktura, optičke mreže i mrežna oprema}
% TODO: Proveriti MACsec i eventualne optičke standarde; podatke proizvođača o akceleraciji označiti kao implementacione tvrdnje, a ne normativne zahteve.

\subsection{Embedded i IoT uređaji}
% TODO: Razmotriti RAM, energiju, latenciju, životni vek ključeva i dostupnost akceleracije; ograničiti broj primera da rad ostane fokusiran.

\subsection{Mapa dokaza za primene}
% TODO: Za svaki primer pripremiti red: standard/verzija, odeljak, AES varijanta, mod, funkcija, obavezna/opciona podrška i datum provere. Izostaviti nepotvrđene primene.
